\chapter{M. N. Srinivas on Caste}

\section{Varna and Jati}

\subsection{From Book View to Field View}

\begin{KeyConcept}
\begin{itemize}
\item \textbf{M. N. Srinivas is the first person to take departure from the book view of Indian society} (the book view given by \textbf{William Jones, Wilkins and Ghurye}) and introduce the \textbf{field view} of Indian society.
\item In the \textbf{book view}, caste becomes \textbf{Varna}; in the \textbf{field view}, caste is \textbf{Jati}. So 'what is caste?' depends on who is speaking, in what context and on which time-scale --- therefore caste cannot be defined easily, and that is why we speak of \textbf{perspectives on caste}.
\item Srinivas was the \textbf{first to discuss the difference between Varna and Jati} (they had been confused before) --- in his article \textbf{'Varna and Jati'}.
\end{itemize}
\end{KeyConcept}

\subsection{Varna vs Jati: The Reference and the Empirical Category}

\begin{KeyConcept}
\begin{itemize}
\item \textbf{Varna} is mentioned in \textbf{Sanskrit texts} (e.g. the \textbf{Rig Veda's Purusha Sukta}), so it is a \textbf{reference category} --- the reference through which we understand the Indian field. The \textbf{Orientalists discovered Varna} by translating texts into English.
\item \textbf{Jati} is \textbf{caste as it actually exists} --- a \textbf{collection of endogamous families inhabiting a region} (local: e.g. the \textbf{Vokkaliga} of \textbf{Rampura}; the \textbf{Thakurs} of UP/Bihar, not Tamil Nadu). Brahmin/Kshatriya are found all over India, but Vokkaliga is not.
\item \textbf{Jati is an empirical category} --- perceivable through the senses (it can be observed); \textbf{Varna is not empirical} --- it is a book/reference category, an \textbf{ideal type} that does not exist in reality.
\end{itemize}
\end{KeyConcept}

\begin{KeyConcept}
Although Varna \textbf{distorts} the picture of caste (showing only four, with Brahmins superior and untouchables inferior all over India --- which is not reality), it has one \textbf{advantage}: it \textbf{enables ordinary men and women to understand and assess the general place of caste throughout India} --- Varna is a \textbf{common language} (a Mishra and a Smarta can recognise each other as Brahmins), bringing stability and order to a chaotic world.
\end{KeyConcept}

\subsection{Four Castes vs Innumerable Jatis}

\begin{KeyConcept}
\begin{itemize}
\item According to the \textbf{Varna scheme there are only four castes} (Brahmin, Kshatriya, Vaishya, Shudra) --- and it \textbf{excludes untouchables} (they are outside the caste structure, the 'broken'/Dalit).
\item But there are \textbf{thousands of Jatis} --- \textbf{innumerable, small and local, with high regional variation}.
\item Example: \textbf{Nama Shudra} is an untouchable jati of West Bengal --- an untouchable is \emph{not} a caste in the Varna model, but is \emph{definitely} a jati.
\item In \textbf{Shripuram} village (Tamil Nadu) the categories are Brahmins, non-Brahmins and \textbf{Adi Dravidas} (untouchables); within the Adi Dravidas are \textbf{Pallans} (Hindu Dalits) and \textbf{Parayans} (Christian Dalits) --- so there is \textbf{hierarchy even within Dalits}, and the notion of purity/impurity exists even among those the Varna model calls impure.
\end{itemize}
\end{KeyConcept}

\subsection{Endogamy, Occupation and the Challenge to Ghurye}

\begin{KeyConcept}
\begin{itemize}
\item In the \textbf{Varna model}, caste is \textbf{endogamous} and \textbf{engaged in a traditional occupation}. In the \textbf{Jati model}, jatis are \textbf{not always endogamous nor always following a traditional occupation}.
\item Here Srinivas is \textbf{challenging Ghurye}: Ghurye's features 'endogamy' and 'lack of unrestricted choice of occupation' are actually the \textbf{Varna (book) view}, not the field view.
\item In his \textbf{Rampura} field study, Srinivas found that \textbf{one jati = one occupation is invalid} --- a single caste group could be involved in multiple occupations (you cannot pin a jati to one occupation).
\item Srinivas is \textbf{not generalising} (unlike Ghurye) --- he is describing Rampura specifically, the field view.
\end{itemize}
\end{KeyConcept}

\subsection{Definite Hierarchy vs Nebulous Hierarchy}

\begin{KeyConcept}
\begin{itemize}
\item In the \textbf{Varna model} there is a \textbf{definite hierarchy and uniformity} --- Brahmins at the top and untouchables at the bottom all over India --- which means \textbf{no mobility}.
\item In the \textbf{Jati model}, hierarchy is \textbf{local and 'nebulous'} (vague, blurred --- a term Srinivas uses). \textbf{Dipankar Gupta} later used the term \textbf{'muddled hierarchy'}.
\item The position of various jatis in the hierarchy is \textbf{not always clear}: e.g. the \textbf{Reddy} claim Kshatriya status (which can never be verified); every jati has its own \textbf{'Jati Puran'} --- a text of mythologies constructed to justify its position (descent from a god, from the moon or sun).
\item The \textbf{Lingayats} (Shaivites who invaded Coorg) claim status equivalent to Brahmins --- unverifiable. The \textbf{Yadavs} are dominant in UP/Bihar but inferior in Maharashtra --- so hierarchy is \textbf{not uniform, not pan-India}.
\end{itemize}
\end{KeyConcept}

\subsection{Why Hierarchy Cannot Be Objectively Assessed}

\begin{KeyConcept}
\begin{itemize}
\item Srinivas held that \textbf{hierarchy cannot be objectively assessed}, for two reasons:
\item 1. \textbf{The problem of criteria} --- should rank be decided by \textbf{numbers, economic position (land ownership), political position, or ritual position}? There is no single deciding factor: some castes are high in numbers but low economically; some are ritually high but economically weak; a poor Brahmin may be ritually superior but economically inferior.
\item 2. \textbf{Claims and counter-claims} --- everyone can claim superiority and everyone can be counter-claimed: in Rampura, the \textbf{Smiths} asserted they would accept cooked (pakka) food only from Brahmins (claiming superiority), while the \textbf{untouchables} asserted they could not accept cooked food and water from Brahmins. So jatis cannot be represented in a simple ladder-like arrangement.
\end{itemize}
\end{KeyConcept}

\subsection{Hierarchy: Manufactured by Anthropologists}

\begin{KeyConcept}
\begin{itemize}
\item Although Srinivas said there is no clear hierarchy at the local level, he \textbf{still placed Brahmins at the top and untouchables at the bottom in Rampura} --- because hierarchy was \textbf{manufactured by anthropologists, not implicit in caste}.
\item As a student of \textbf{social anthropology} he was a \textbf{structural functionalist} (a follower of \textbf{Evans-Pritchard and Radcliffe-Brown}), committed to showing how parts make up the system --- so he constructed hierarchy to show how it maintains the caste system.
\end{itemize}
\end{KeyConcept}

\subsection{Ranking: Ritual vs Secular Factors}

\begin{KeyConcept}
\begin{itemize}
\item In the \textbf{Varna model}, the criterion of ranking is \textbf{ritual} --- the notion of \textbf{purity and pollution} --- stressing \textbf{attributional features} (attributes/traits: food, skin colour, dress, language, mode of production).
\item In the \textbf{Jati model}, rank is shaped by \textbf{political and economic resources} --- the \textbf{secular factors} --- the same criteria as for \textbf{dominance} (the \textbf{dominant caste}, originally called 'dominant jati').
\item Examples from Rampura: the \textbf{Vokkaliga} (peasants) are Shudras in the Varna model (inferior), but in the field they are \textbf{dominant} --- prosperous, wealthy, landowners, and numerous (the OBCs have the maximum population in India) --- claiming superiority even over Brahmins.
\item \textbf{Not all Brahmins are pure} --- in the field view one may find impure/polluted Brahmins.
\item \textbf{Mobility}: in the Varna model mobility comes only through \textbf{ritual considerations} (anyone can become Brahmin/vegetarian); in the Jati model mobility comes through \textbf{numbers, economic position and political position} --- rituals are not the only route.
\end{itemize}
\end{KeyConcept}

\begin{KeyConcept}
Final remark: there is a \textbf{disparity between ritual and secular position} --- a caste can be \textbf{ritually high but ranked lower} in the local hierarchy, which is determined by \textbf{secular factors} (economic, political and educational status).
\end{KeyConcept}

\subsection{Dominant Caste and the Role of Context}

\begin{KeyConcept}
\begin{itemize}
\item Asked who the dominant caste of Rampura is, Srinivas answered: \textbf{'it is the context that defines your status.'}
\item In a \textbf{secular context} (the village headman is a Vokkaliga, and Brahmins depend on the Vokkaliga headman), economic and political factors make the \textbf{Vokkaliga} dominant.
\item In a \textbf{marriage/ritual context} (the Vokkaliga depend on Brahmins), the \textbf{Brahmins} are dominant.
\item So there is \textbf{no universal idea of dominance}: the \emph{concept} of dominant caste is about \textbf{secular rank} (political, economic, numerical factors), but actual dominance depends on context.
\end{itemize}
\end{KeyConcept}

\begin{KeyConcept}
\begin{itemize}
\item \textbf{Caste is immutable} (fixed) in the \textbf{Varna model} (no mobility), but \textbf{mutable} in the \textbf{Jati model} --- it is in the field that Srinivas observed \textbf{Sanskritization} (social mobility).
\item \textbf{Census is a source of social mobility} (per Srinivas): ten years back you were Kshatriya, now you claim Brahmin status --- these claims can never be objectively assessed.
\item Castes construct \textbf{mythologies and texts} (Jati Purans) to claim ritual superiority, even after having secured secular power (numbers, wealth, political power).
\end{itemize}
\end{KeyConcept}

\begin{QuickRevision}
\begin{itemize}
\item Srinivas: first to depart from book view (William Jones/Wilkins/Ghurye) to field view; Varna = book, Jati = field; article 'Varna and Jati'.
\item Varna = reference/ideal type (Sanskrit texts, Orientalists), not empirical; Jati = empirical (observable), 'collection of endogamous families inhabiting a region'.
\item Varna = four castes (excludes untouchables); Jati = innumerable/local (Nama Shudra; Shripuram's Adi Dravidas, Pallans vs Parayans).
\item Varna: endogamy + traditional occupation; Jati: not always endogamous, one jati $\neq$ one occupation (Rampura) --- challenges Ghurye.
\item Varna: definite hierarchy/uniformity/no mobility; Jati: local, 'nebulous' hierarchy (Dipankar Gupta's 'muddled'); Reddy/Jati Puran/Lingayat/Yadav examples.
\item Hierarchy not objectively assessable: criteria problem + claims/counter-claims (Smiths vs untouchables of Rampura); hierarchy manufactured by anthropologists (structural functionalism).
\item Ranking: ritual (purity/pollution, attributional) vs secular (political/economic); Vokkaliga dominant; not all Brahmins pure; mobility via secular factors too.
\item Dominant caste = secular rank, but 'context defines status' (secular vs ritual context); caste immutable (Varna) vs mutable (Jati, Sanskritization); census = mobility.
\end{itemize}
\end{QuickRevision}

\section{Features of Jati, Dominant Caste and Vote Bank}

\subsection{Varna  to  Caste  to  Sub-Caste (Fission)}

\begin{KeyConcept}
\begin{itemize}
\item After distinguishing \textbf{Varna and Jati}, Srinivas discusses the features of caste --- but he says \textbf{features of Jati}, not caste (he is a field-view scholar; the English for jati in the caste context is \textbf{sub-caste}).
\item The structure he works with: \textbf{Varna $\rightarrow$ Caste $\rightarrow$ Sub-caste}. Within Brahmin there are multiple castes (\textbf{Sri Vaishnava, Smarta, Bhatkal}, etc.), and within each caste there are multiple sub-castes; within Shudra there are multiple castes, each divided into sub-castes.
\item This is \textbf{fission} --- and Srinivas is largely studying the fission process (as against fusion).
\end{itemize}
\end{KeyConcept}

\subsection{Features of the Sub-Caste (Jati)}

\begin{KeyConcept}
\begin{itemize}
\item Srinivas's features of jati/sub-caste (largely from his field work, in contrast to Ghurye's book-view features):
\item 1. \textbf{Hierarchy} --- but \emph{nebulous}.
\item 2. \textbf{Purity and pollution} --- but with no clear binary.
\item 3. \textbf{Occupational association} --- systematization of occupational differentiation.
\item 4. \textbf{Restrictions on commensality, dress, speech and custom}.
\item 5. \textbf{Caste panchayat and assemblies}.
\item 6. \textbf{Sub-caste as a unit of endogamy}.
\item 7. \textbf{Members share a common culture} (each sub-caste has its own distinct social customs and ritual life).
\end{itemize}
\end{KeyConcept}

\begin{ComparisonBox}
\begin{itemize}
\item Ghurye discussed features of \textbf{caste} (via the \textbf{Varna/book view} --- definite hierarchy); Srinivas discusses features of \textbf{jati/sub-caste} (via the \textbf{field view} --- nebulous hierarchy).
\item Srinivas witnessed a \textbf{completely different picture of jati in the field, in opposition to Ghurye's analysis based on ancient Sanskritic texts}.
\end{itemize}
\end{ComparisonBox}

\subsection{Hierarchy: Nebulous}

\begin{KeyConcept}
\begin{itemize}
\item In the book view, Brahmins are at the top and untouchables at the bottom (definite hierarchy). In the field, the \textbf{middle ranks are highly nebulous} (vague) --- we cannot make an objective claim about who is where.
\item There are \textbf{claims and counter-claims}: many untouchables do not even accept water or food from the hands of Brahmins.
\item Yet, despite these contradictions, Srinivas \textbf{recorded Brahmins at the top and untouchables at the bottom even in the field} --- because of his \textbf{commitment to structural functionalism}: a system must have a definite hierarchy to be called a system. This commitment made him \textbf{oblivious to the inconsistencies between book view and field view}.
\end{itemize}
\end{KeyConcept}

\subsection{Purity and Pollution: No Clear Binary}

\begin{KeyConcept}
\begin{itemize}
\item Srinivas saw that Brahmins are considered the most pure and untouchables the most polluted --- but in the field he also saw \textbf{polluted people claiming purity}.
\item Example: the \textbf{Smiths} (untouchables) call themselves \textbf{'Vishwakarma Brahmins'}.
\item There is \textbf{no clear-cut binary distinction between pure and impure} in the field. What marks someone as pure/impure is only \textbf{ritual considerations} --- and even the low castes (or those outside the caste) can imbibe those rituals through \textbf{Sanskritization}.
\end{itemize}
\end{KeyConcept}

\subsection{Occupational Association and Nesfield}

\begin{KeyConcept}
\begin{itemize}
\item \textbf{Jati is the systematization of occupational differentiation} --- jatis are occupational groups, known by their occupations (\textbf{Sonar} = goldsmith, \textbf{Lohar} = ironsmith, \textbf{Kumhar} = potter, \textbf{Chamar} = leather worker).
\item Srinivas drew on \textbf{Nesfield's occupational theory of caste} (not Risley's racial theory).
\item But \textbf{one jati $\neq$ one occupation as a universal rule} --- in the field a jati may have multiple occupations (someone from another jati working a non-traditional occupation).
\end{itemize}
\end{KeyConcept}

\subsection{Restrictions, Caste Panchayat, Endogamy, Common Culture}

\begin{KeyConcept}
\begin{itemize}
\item \textbf{Restrictions on commensality, dress, speech and custom} --- even in the field, the lower castes do not accept water/food from the upper castes (and claim Brahminical status).
\item \textbf{Caste panchayat and assemblies} --- their function is to ensure \textbf{compliance with caste norms} (no one digresses/disrupts the norms), so they maintain the system as it is and prevent intermixing.
\item \textbf{Sub-caste is a unit of endogamy}, and \textbf{members share a common culture} (distinct social customs and ritual life within the same caste).
\end{itemize}
\end{KeyConcept}

\subsection{The Dominant Jati (Dominant Caste)}

\begin{KeyConcept}
\begin{itemize}
\item \textbf{Dominant caste} is a concept of the \textbf{jati (field) model} --- Srinivas does not say Brahmins are dominant; he points to the \textbf{Vokkaliga} of Rampura.
\item In the \textbf{ritual hierarchy} (Sanskrit texts) Brahmins are dominant, but in the field someone else is dominant --- and \textbf{context defines status}.
\item The \textbf{three main factors of dominance}: (1) \textbf{numerical preponderance}, (2) \textbf{economic ownership of land}, (3) \textbf{political position}.
\end{itemize}
\end{KeyConcept}

\begin{ExampleBox}
In \textbf{Rampura}, the \textbf{Vokkaliga} are dominant: they are peasants (i.e. Shudra varna), but they have numerical preponderance, ownership of land, accumulated wealth, and the village \textbf{headman} is a Vokkaliga (political position).
\end{ExampleBox}

\subsection{Westernization and the New Factors of Dominance}

\begin{KeyConcept}
\begin{itemize}
\item With \textbf{Westernization}, the factors of dominance recorded \textbf{more entries}:
\item 1. \textbf{Representation in government jobs/bureaucracy} (IAS/IPS, even income-tax inspector).
\item 2. \textbf{Urban income} (business).
\item 3. \textbf{English-medium education}.
\item 4. \textbf{Ritual rank} --- still important: a Shudra who owns land, has numbers and political power will claim high status even at the ritual level (the 'final nail in the coffin').
\end{itemize}
\end{KeyConcept}

\begin{ExampleBox}
\begin{itemize}
\item After \textbf{land reforms} (the land ceiling/land sealing acts and land fragmentation), the Vokkaliga lost their hold over land --- so they began sending their children to \textbf{English-medium schools} to retain dominance.
\item With \textbf{reservation}, ritually low people became bureaucrats and acquired land --- their say in the village became dominant even without numerical preponderance.
\item So \textbf{there has been a shift in the idea of dominance since independence} --- earlier numerical preponderance, headman and peasant status sufficed; today Westernization and reservation have added new factors.
\end{itemize}
\end{ExampleBox}

\subsection{Vote Bank and Horizontal/Vertical Solidarity}

\begin{KeyConcept}
\begin{itemize}
\item \textbf{Vote bank} is the \textbf{modern avatar (20th century) of caste}. Srinivas recorded it in his book \textbf{'Caste: Its 20th Century Avatar'} --- caste as a political category in modern times, the \textbf{structural differentiation} (Parsons) of caste.
\item The first person to talk about the \textbf{modern avatar of caste was Ghurye} (caste patriotism, caste association); Srinivas called caste associations \textbf{horizontal solidarity}.
\item A vote bank is the \textbf{fusion of castes} (multiple castes fusing to form a party, in Weber's language). Example: the \textbf{Yadavs of UP} were once a vote bank, but now are spread across SP, Congress, BJP and BSP --- so there is no neat caste vote bank today.
\end{itemize}
\end{KeyConcept}

\begin{KeyConcept}
\begin{itemize}
\item Srinivas's two concepts:
\item \textbf{Horizontal solidarity} --- castes that \textbf{cannot be ranked hierarchically} (jatis cannot be ranked) form associations; this happens at the \textbf{local level} (nebulous hierarchy).
\item \textbf{Vertical solidarity} --- upper and lower castes integrated for \textbf{common interests} (e.g. the appointment of a Dalit President was an attempt by the ruling party to establish ties with the lower caste and shed its 'Brahminical/anti-lower-caste' image, forming a vote bank). At the \textbf{national level} the Varna model operates, producing vertical solidarity.
\item These two terms are inspired by / linked to \textbf{Rudolph and Rudolph} (Lloyd and Susanne Rudolph, Padma Bhushan-awarded political scientists), who gave \textbf{vertical mobilization, horizontal mobilization and differential mobilization}.
\end{itemize}
\end{KeyConcept}

\begin{ExampleBox}
Answer tip: in a question on \textbf{Durkheim's mechanical/organic solidarity}, adding Srinivas's \textbf{horizontal and vertical solidarity} at the end fetches good marks; in a question on \textbf{Marx's historical materialism}, adding that A. R. Desai saw 'bourgeois industrialization instead of socialism' (the Marxian mode-of-production theory failing in India) does the same.
\end{ExampleBox}

\subsection{Conclusion}

\begin{KeyConcept}
\begin{itemize}
\item In conclusion, \textbf{Srinivas's perspective helped the sociology of caste become the sociology of jati --- an empirical category}, unlike the highly abstract 'caste'.
\item Clarification of terms: \textbf{Varna} is discovered by Orientalists (found in texts); \textbf{Jati} is found in reality; \textbf{caste} is a category made by colonialists/census commissioners. India had \textbf{only Varna and Jati} --- 'caste' is either the Varna model or the Jati model.
\end{itemize}
\end{KeyConcept}

\begin{QuickRevision}
\begin{itemize}
\item Srinivas discusses features of \textbf{jati/sub-caste} (field view), not caste: Varna $\rightarrow$ Caste $\rightarrow$ Sub-caste (fission).
\item Features of jati: (nebulous) hierarchy; purity/pollution (no binary); occupational association (Nesfield, not Risley); restrictions on commensality/dress/speech/custom; caste panchayat; sub-caste as endogamous unit; common culture.
\item Nebulous hierarchy --- middle ranks vague; but he still placed Brahmins top/untouchables bottom (commitment to structural functionalism).
\item Smiths = 'Vishwakarma Brahmins' (polluted claiming purity); purity/pollution = ritual, mutable via Sanskritization.
\item Dominant caste = jati model; three factors (numerical preponderance, land ownership, political position) + Westernization's new factors (government jobs, urban income, English education, ritual rank); Rampura's Vokkaliga; land reforms $\rightarrow$ English-medium shift; reservation $\rightarrow$ new dominance.
\item Vote bank = modern avatar (20th century) of caste; Ghurye's caste patriotism $\rightarrow$ Srinivas's horizontal solidarity; fusion of castes; structural differentiation.
\item Horizontal solidarity (unranked, local) vs vertical solidarity (upper+lower for common interest, national/Varna); Rudolph \& Rudolph's vertical/horizontal/differential mobilization.
\item Conclusion: sociology of caste became sociology of jati (empirical).
\end{itemize}
\end{QuickRevision}

